\clearpage\section{Does more Gaussian entanglement imply a stronger clock?}
\subsection{Hold the quadratic state fixed and vary the cubic vertex}
At fixed $(M,g)$, the quadratic state fixes the reduced covariance $V_q$ of the light field and its conjugate momentum. With $\hbar=1$, $\nu_q=\sqrt{\det V_q}$ and $n=\nu_q-1/2$. The entropy per real sector is
\[
 S_q=(n+1)\ln(n+1)-n\ln n .
\]
We evaluate this Gaussian entropy at $x=0.01$. The quadratic Hamiltonian contains $M$ and $g$. Changing $\lambda_3$ leaves its covariance calculation unchanged. By contrast the leading bispectrum contains exactly one cubic vertex. Fixing $r$, the two final powers and the three mixed propagators therefore gives
\[
 F(r;\lambda_3)=\frac{\lambda_3}{\lambda_{3,\rm ref}}F(r;\lambda_{3,\rm ref}) .
\]
At $(M,g)=(2,0.1)$ the reference is $\lambda_{3,\rm ref}/H=0.01$. The last column below uses the oscillatory part from the full fit. For a known template $T_i$ at 50 ratios and independent common absolute noise $s=10^{-10}$, the diagnostic is $[\sum_i(T_i/s)^2]^{1/2}$. Scaling the template scales this value by the same absolute factor.
\begin{center}\begin{tabular}{rrrr}\toprule $\lambda_3/H$ & $S_q$ (nats) & $F(0.03)$ & Known-template S/N\\\midrule
0 & 0.000184599 & 0 & 0\\
0.005 & 0.000184599 & $-1.99291\times10^{-10}$ & 49.81\\
0.010 & 0.000184599 & $-3.98581\times10^{-10}$ & 99.63\\\bottomrule\end{tabular}\end{center}
The zero-cubic row gives a direct counterexample to the claim that nonzero Gaussian entanglement guarantees a leading clock signal. The assumed noise and known waveform define a controlled comparison. An observational signal-to-noise forecast would also require a survey model and covariance.
\fig{entropy_signal_control}{.9}{Left: varying quadratic mass or mixing changes the entropy and also changes the signed shape at one triangle. Bars show regulator sensitivity. Right: varying only the cubic coupling scales the shape while holding Gaussian entropy fixed.}

\clearpage\subsection{Changing mass or mixing probes a different question}
With $\lambda_3/H=0.01$ and $r=0.03$, the selected five models yield the following results. The power column is $100(P_q/P_q^{(0)}-1)$ at $x=0.01$. It compares the light power with the free value at the same time.
\begin{center}\begin{tabular}{rrrrr}\toprule $M$ & $g$ & $10^3 S_q$ (nats) & Power change (\%) & $10^{10}F(0.03)$\\\midrule
2.0 & 0.05 & 0.05094 & 0.09297 & $-0.5084$\\
2.0 & 0.10 & 0.18460 & 0.37140 & $-3.9858$\\
2.0 & 0.20 & 0.67903 & 1.47821 & $-29.3588$\\
1.8 & 0.10 & 1.37718 & 0.57652 & $+0.3916$\\
2.4 & 0.10 & 0.07630 & 0.18252 & $+0.4488$\\\bottomrule\end{tabular}\end{center}
Along the fixed-$M=2$ slice, increasing $g$ increases the entropy and the magnitude at this triangle. Comparing $M=1.8$ with $M=2$ at the same $g=0.1$ reverses that ordering: the lighter case has higher entropy and a smaller $|F(0.03)|$. Changing the mass shifts the clock phase, so the sign and size at a single triangle cannot determine the multi-triangle clock amplitude. A full fit like the one in Section 10 would be needed at each mass to compare fitted amplitudes.
\source{10}{Parameter comparison, cubic scaling and common-noise template control}

\subsection{Conclusion}
The quadratic mixing transfers fluctuations between the two fields. Tracing out the spectator leaves Gaussian entanglement that can be measured from the light covariance. The cubic interaction probes more information: its leading three-point term integrates the relative phases of three unequal-time mixed propagators. At fixed covariance its strength can be varied continuously through zero.

For the baseline with a nonzero cubic coupling, the numerical time-integral errors lie below one percent of the fitted clock envelope. The spread $m/H=1.993$--$2.020$ comes from fitting roughly one cycle across different backgrounds and momentum windows. Thus the model can encode the heavy mass in the squeezed shape, while Gaussian entropy alone fixes neither the size of that signal nor how accurately its frequency can be recovered.

\subsection*{References for the bispectrum calculation}
\begin{enumerate}\small\setlength{\itemsep}{1pt}
\item S. Weinberg, \emph{Quantum Contributions to Cosmological Correlations}, Phys. Rev. D 72, 043514 (2005). \href{https://arxiv.org/abs/hep-th/0506236}{arXiv:hep-th/0506236}.
\item X. Chen and Y. Wang, \emph{Quasi-Single Field Inflation and Non-Gaussianities}, JCAP 04, 027 (2010). \href{https://arxiv.org/abs/0911.3380}{arXiv:0911.3380}.
\item N. Arkani-Hamed and J. Maldacena, \emph{Cosmological Collider Physics} (2015). \href{https://arxiv.org/abs/1503.08043}{arXiv:1503.08043}.
\end{enumerate}
